% GENERATED FILE. Edit canonical lessons, then run npm run generate:books.
% canonical-source-hash: fnv1a64:4beb9ae0a3af34a2
% canonical-lessons: ML-C554-sahodari, ML-C554-matapitakkal, ML-C554-aliyan, ML-C554-nattun, ML-C554-kutumbamgam

\chapter{Parents and In-laws}
\label{ch:ml-parents-and-in-laws}
\hlchaptermodality{554}

\begin{chapteropening}
\textbf{By the end of this chapter:} \emph{I can say a sister, parents, a brother-in-law, a sister-in-law and a family member.}

Complete the last lesson of chapter 554: I can say a sister, parents, a brother-in-law, a sister-in-law and a family member.
\end{chapteropening}

\section[\texorpdfstring{\ml{സഹോദരി} (sahōdari)}{സഹോദരി (sahōdari)}]{\ml{സഹോദരി} (sahōdari) — a sister}

\label{lesson:ML-C554-sahodari}



\begin{quote}
Before the new one: say the Malayalam for a sibling, then the Malayalam for a brother.
\end{quote}

\subsection*{You'll want to know: \ml{സഹോദരി}}
\textbf{\ml{സഹോദരി}} — \emph{sahōdari} — \textquotedblleft{}a sister\textquotedblright{}.

\begin{quote}
\ml{അവന്റെ} \ml{സഹോദരി} \ml{ഒരു} \ml{നഴ്സാണ്}.
\end{quote}

\begin{quote}
\emph{avanṟe sahōdari oru naḻsāṇŭ.}
\end{quote}

\begin{quote}
His sister is a nurse.
\end{quote}

At home, \ml{ചേച്ചി} or \ml{അനിയത്തി}.

\begin{cousinweb}
\textbf{\ml{സഹോദരി}} (\emph{sahōdari}) is borrowed from Sanskrit \emph{sahodarī}, the feminine of \emph{sahodara}, \textquotedblleft{}born of the same womb\textquotedblright{} (Monier-Williams).
\end{cousinweb}

\begin{grammarlens}[title={What you've built}]
One more word for describing what you lived and felt.
\end{grammarlens}

\subsection*{Your turn}
\begin{itemize}
\raggedright
  \item \emph{Say it:} \emph{sahōdari}
  \item \emph{Say it:} \emph{sahōdari}, once more
  \item \emph{Say it:} \emph{avanṟe sahōdari oru naḻsāṇŭ.} — His sister is a nurse.
  \item \emph{Recall:} say the Malayalam for a sibling, then the Malayalam for a brother, then say \emph{sahōdari} again
\end{itemize}

\begin{tcolorbox}[breakable,colback=teal!4,colframe=teal!35!black,title={Before you move on}]
What does \ml{സഹോദരി} mean? (\textquotedblleft{}A sister\textquotedblright{}.) Say it once more.
\end{tcolorbox}
\section[\texorpdfstring{\ml{മാതാപിതാക്കൾ} (mātāpitākkaḷ)}{മാതാപിതാക്കൾ (mātāpitākkaḷ)}]{\ml{മാതാപിതാക്കൾ} (mātāpitākkaḷ) — parents}

\label{lesson:ML-C554-matapitakkal}



\begin{quote}
Before the new one: say the Malayalam for a brother, then the Malayalam for a sister.
\end{quote}

\subsection*{You'll want to know: \ml{മാതാപിതാക്കൾ}}
\textbf{\ml{മാതാപിതാക്കൾ}} — \emph{mātāpitākkaḷ} — \textquotedblleft{}parents\textquotedblright{}.

\begin{quote}
\ml{എന്റെ} \ml{മാതാപിതാക്കൾ} \ml{ഗ്രാമത്തിലാണ്} \ml{താമസിക്കുന്നത്}.
\end{quote}

\begin{quote}
\emph{enṟe mātāpitākkaḷ grāmattilāṇŭ tāmasikkunnatŭ.}
\end{quote}

\begin{quote}
My parents live in the village.
\end{quote}

\ml{മാതാവ്} mother + \ml{പിതാവ്} father, with the plural -\ml{ക്കൾ}.

\begin{cousinweb}
\textbf{\ml{മാതാപിതാക്കൾ}} (\emph{mātāpitākkaḷ}) is a Malayalam compound built from Sanskrit \emph{mātṛ}, \textquotedblleft{}mother\textquotedblright{}, and \emph{pitṛ}, \textquotedblleft{}father\textquotedblright{} (Monier-Williams), with the plural ending -\ml{ക്കൾ}.
\end{cousinweb}

\begin{grammarlens}[title={What you've built}]
One more word for describing what you lived and felt.
\end{grammarlens}

\subsection*{Your turn}
\begin{itemize}
\raggedright
  \item \emph{Say it:} \emph{mātāpitākkaḷ}
  \item \emph{Say it:} \emph{mātāpitākkaḷ}, once more
  \item \emph{Say it:} \emph{enṟe mātāpitākkaḷ grāmattilāṇŭ tāmasikkunnatŭ.} — My parents live in the village.
  \item \emph{Recall:} say the Malayalam for a brother, then the Malayalam for a sister, then say \emph{mātāpitākkaḷ} again
\end{itemize}

\begin{tcolorbox}[breakable,colback=teal!4,colframe=teal!35!black,title={Before you move on}]
What does \ml{മാതാപിതാക്കൾ} mean? (\textquotedblleft{}Parents\textquotedblright{}.) Say it once more.
\end{tcolorbox}
\section[\texorpdfstring{\ml{അളിയൻ} (aḷiyan)}{അളിയൻ (aḷiyan)}]{\ml{അളിയൻ} (aḷiyan) — a brother-in-law}

\label{lesson:ML-C554-aliyan}



\begin{quote}
Before the new one: say the Malayalam for a sister, then the Malayalam for parents.
\end{quote}

\subsection*{You'll want to know: \ml{അളിയൻ}}
\textbf{\ml{അളിയൻ}} — \emph{aḷiyan} — \textquotedblleft{}a brother-in-law\textquotedblright{}.

\begin{quote}
\ml{എന്റെ} \ml{അളിയൻ} \ml{ഒരു} \ml{എഞ്ചിനീയറാണ്}.
\end{quote}

\begin{quote}
\emph{enṟe aḷiyan oru eñcinīyaṟāṇŭ.}
\end{quote}

\begin{quote}
My brother-in-law is an engineer.
\end{quote}

A wife's brother or a sister's husband.

\begin{cousinweb}
\textbf{\ml{അളിയൻ}} (\emph{aḷiyan}) is native Dravidian: DEDR 2410 gives it as a wife's brother or a sister's husband (DEDR 2410).
\end{cousinweb}

\begin{grammarlens}[title={What you've built}]
One more word for describing what you lived and felt.
\end{grammarlens}

\subsection*{Your turn}
\begin{itemize}
\raggedright
  \item \emph{Say it:} \emph{aḷiyan}
  \item \emph{Say it:} \emph{aḷiyan}, once more
  \item \emph{Say it:} \emph{enṟe aḷiyan oru eñcinīyaṟāṇŭ.} — My brother-in-law is an engineer.
  \item \emph{Recall:} say the Malayalam for a sister, then the Malayalam for parents, then say \emph{aḷiyan} again
\end{itemize}

\begin{tcolorbox}[breakable,colback=teal!4,colframe=teal!35!black,title={Before you move on}]
What does \ml{അളിയൻ} mean? (\textquotedblleft{}A brother-in-law\textquotedblright{}.) Say it once more.
\end{tcolorbox}
\section[\texorpdfstring{\ml{നാത്തൂൻ} (nāttūn)}{നാത്തൂൻ (nāttūn)}]{\ml{നാത്തൂൻ} (nāttūn) — a sister-in-law}

\label{lesson:ML-C554-nattun}



\begin{quote}
Before the new one: say the Malayalam for parents, then the Malayalam for a brother-in-law.
\end{quote}

\subsection*{You'll want to know: \ml{നാത്തൂൻ}}
\textbf{\ml{നാത്തൂൻ}} — \emph{nāttūn} — \textquotedblleft{}a sister-in-law\textquotedblright{}.

\begin{quote}
\ml{നാത്തൂനും} \ml{ഞാനും} \ml{ഒരുമിച്ച്} \ml{ചന്തയിൽ} \ml{പോയി}.
\end{quote}

\begin{quote}
\emph{nāttūnuṁ ñānuṁ orumiccŭ cantayil pōyi.}
\end{quote}

\begin{quote}
My sister-in-law and I went to the market together.
\end{quote}

A husband's sister or a brother's wife.

\begin{cousinweb}
\textbf{\ml{നാത്തൂൻ}} (\emph{nāttūn}) is native Dravidian: DEDR 3644 gives it as a sister-in-law, a brother's wife (DEDR 3644).
\end{cousinweb}

\begin{grammarlens}[title={What you've built}]
One more word for describing what you lived and felt.
\end{grammarlens}

\subsection*{Your turn}
\begin{itemize}
\raggedright
  \item \emph{Say it:} \emph{nāttūn}
  \item \emph{Say it:} \emph{nāttūn}, once more
  \item \emph{Say it:} \emph{nāttūnuṁ ñānuṁ orumiccŭ cantayil pōyi.} — My sister-in-law and I went to the market together.
  \item \emph{Recall:} say the Malayalam for parents, then the Malayalam for a brother-in-law, then say \emph{nāttūn} again
\end{itemize}

\begin{tcolorbox}[breakable,colback=teal!4,colframe=teal!35!black,title={Before you move on}]
What does \ml{നാത്തൂൻ} mean? (\textquotedblleft{}A sister-in-law\textquotedblright{}.) Say it once more.
\end{tcolorbox}
\section[\texorpdfstring{\ml{കുടുംബാംഗം} (kuṭumbāṁgaṁ)}{കുടുംബാംഗം (kuṭumbāṁgaṁ)}]{\ml{കുടുംബാംഗം} (kuṭumbāṁgaṁ) — a family member}

\label{lesson:ML-C554-kutumbamgam}



\begin{quote}
Before the new one: say the Malayalam for a brother-in-law, then the Malayalam for a sister-in-law.
\end{quote}

\subsection*{You'll want to know: \ml{കുടുംബാംഗം}}
\textbf{\ml{കുടുംബാംഗം}} — \emph{kuṭumbāṁgaṁ} — \textquotedblleft{}a family member\textquotedblright{}.

\begin{quote}
\ml{കല്യാണത്തിന്} \ml{എല്ലാ} \ml{കുടുംബാംഗങ്ങളും} \ml{എത്തി}.
\end{quote}

\begin{quote}
\emph{kalyāṇattinŭ ellā kuṭumbāṁgaṅṅaḷuṁ etti.}
\end{quote}

\begin{quote}
All the family members came to the wedding.
\end{quote}

Plural: \ml{കുടുംബാംഗങ്ങൾ}, \emph{kuṭumbāṁgaṅṅaḷ}.

\ml{കുടുംബം} family + \ml{അംഗം} a limb, a member.

\begin{cousinweb}
\textbf{\ml{കുടുംബാംഗം}} (\emph{kuṭumbāṁgaṁ}) is a learned compound built from Sanskrit \emph{kuṭumba}, \textquotedblleft{}a household\textquotedblright{} (CDIAL 3233), and \emph{aṅga}, \textquotedblleft{}a limb\textquotedblright{} (CDIAL 114).
\end{cousinweb}

\begin{grammarlens}[title={What you've built}]
One more word for describing what you lived and felt.
\end{grammarlens}

\subsection*{Your turn}
\begin{itemize}
\raggedright
  \item \emph{Say it:} \emph{kuṭumbāṁgaṁ}
  \item \emph{Say it:} \emph{kuṭumbāṁgaṁ}, once more
  \item \emph{Say it:} \emph{kalyāṇattinŭ ellā kuṭumbāṁgaṅṅaḷuṁ etti.} — All the family members came to the wedding.
  \item \emph{Recall:} say the Malayalam for a brother-in-law, then the Malayalam for a sister-in-law, then say \emph{kuṭumbāṁgaṁ} again
\end{itemize}

\begin{tcolorbox}[breakable,colback=teal!4,colframe=teal!35!black,title={Before you move on}]
What does \ml{കുടുംബാംഗം} mean? (\textquotedblleft{}A family member\textquotedblright{}.) Say it once more.
\end{tcolorbox}
